\expandafter\def\csname HistoricalNoteText21\endcsname{On 4 Feb. elements III give a right ascension smaller by $13'$, and on 28 June one smaller by $32'$. They give a declination greater by $2'$.}
\SourceLeaf{228}%
\bigskip\Needspace{10\baselineskip}
\pdfbookmark[1]{Gauss to Bode. 1802 Sept. 10}{bode18020910}
\begin{center}GAUSS TO BODE.\end{center}
\begin{flushright}Braunschweig 1802 Sept.\ 10.\end{flushright}
\begin{center}
\medskip\hrule\medskip
Astronomisches Jahrbuch for 1805. Berlin 1802. Page 227.\par
\medskip\hrule
\end{center}\medskip

The elements used to calculate the following ephemeris for the 9th primary planet, \textit{Pallas Olbersiana}, are based on the earliest meridian observations at Seeberg and the latest at Greenwich, and on Dr.\ \textsc{Olbers}’s observations of 8 July. They are given below.
\begin{center}
{\small\setlength{\tabcolsep}{3pt}
\begin{tabular}{l r}
Epoch 1802 March 31, Seeberg meridian $\ .\ .\ .\ .\ .$&$162^0\ 55'\quad6''8$\\
Daily tropical motion $\ .\ .\ .\ .\ .\ .\ .\ .\ .$&$769''7263$\\
Logarithm of the semimajor axis $\ .\ .\ .\ .\ .\ .$&$0.4424992$\\
Aphelion $\ .\ .\ .\ .\ .\ .\ .\ .\ .\ .\ .\ .\ .\ .\ .$&$301^0\ 38'\ 42''$\\
Eccentricity $\ .\ .\ .\ .\ .\ .\ .\ .\ .\ .\ .\ .\ .\ .\ .$&$0.243888$\\
Node {\AstroSymbols ☊} $\ .\ .\ .\ .\ .\ .\ .\ .\ .\ .\ .\ .\ .\ .\ .\ .$&$172^0\ 26'\ 31''$\\
Inclination $\ .\ .\ .\ .\ .\ .\ .\ .\ .\ .\ .\ .\ .\ .\ .\ .\ .$&$34\ 36\ 59$
\end{tabular}}
\end{center}
If we set the brightness of the planet on 8 July of this year $=1$, it was as follows.
\begin{center}
\begin{tabular}{r l r}
on 4 April&$\ .\ .\ .\ .\ .\ .\ .\ .\ .\ .$&$=4.23$\\
4 May&$\ .\ .\ .\ .\ .\ .\ .\ .\ .\ .$&$=2.75$\\
26 July&$\ .\ .\ .\ .\ .\ .\ .\ .\ .\ .$&$=0.80$\\
24 August&$\ .\ .\ .\ .\ .\ .\ .\ .\ .$&$=0.60$\\
\multicolumn{2}{l}{greatest brightness next year}&$=0.653$
\end{tabular}
\end{center}

This shows that we cannot yet be quite sure of finding the planet again in 1803. Its great height above the horizon and its greater distance from the twilight will have to make up for the loss. Still, I flatter myself that it will be easier to find in southern climates, for example in Palermo, where the June nights are completely dark at midnight.

If I later calculate an ephemeris of $\CeresSymbol$ from improved elements, I shall be sure to send it to Your — as well\HistoricalNote{21}. [The ephemeris covers Feb.\ 4 to June 28 at midnight on the Seeberg meridian in 1803. It gives R.A. and declination to the nearest minute at intervals of three days.]
\begin{center}
\begin{tabular}{l}
{\AstroSymbols ☍} June 30. Afternoon.\\
according to elements III, June 30. Morning, civil time.
\end{tabular}
\end{center}
